\documentclass{article}
\usepackage{amsmath,amsthm,amssymb}
\usepackage{algorithm,algorithmic}
\usepackage{bm}
\usepackage{enumitem}
\newtheorem{theorem}{Theorem}
\newtheorem{lemma}{Lemma}
\newtheorem{assumption}{Assumption}
\newtheorem{corollary}{Corollary}
\newcommand{\E}{\mathbb{E}}
\newcommand{\R}{\mathbb{R}}
\newcommand{\norm}[1]{\left\lVert #1\right\rVert}
\newcommand{\grad}{\nabla}
\begin{document}

\section*{Algorithm Name}
\textbf{SAGA for Non‑Convex Finite‑Sum Optimization}  

We consider the finite‑sum problem  

\[
\min_{x\in\R^{d}} \; f(x)\;:=\;\frac{1}{n}\sum_{i=1}^{n} f_i(x),
\]

where each component $f_i$ is possibly non‑convex but $L$‑smooth.

--------------------------------------------------------------------

\section*{Mathematical Formulation}

\begin{itemize}
  \item \textbf{Parameters:}
    \begin{itemize}
      \item step size $\eta>0$,
      \item number of component functions $n$,
      \item initial point $x_{0}\in\R^{d}$,
      \item table of past points $\{\phi_i\}_{i=1}^{n}$, initialised as $\phi_i=x_{0}$.
    \end{itemize}
  \item \textbf{Iteration $k=0,1,\dots$:}
    \begin{enumerate}[label=\arabic*.]
      \item Sample an index $i_k$ uniformly from $\{1,\dots ,n\}$.
      \item Compute the \emph{variance‑reduced gradient estimator}
      \[
      g_k \;=\; \grad f_{i_k}(x_k)\;-\;\grad f_{i_k}(\phi_{i_k})\;+\;\frac{1}{n}\sum_{j=1}^{n}\grad f_{j}(\phi_{j}).
      \tag{1}
      \]
      \item Update the iterate
      \[
      x_{k+1}=x_k-\eta\, g_k . \tag{2}
      \]
      \item Update the table entry of the sampled component
      \[
      \phi_{i_k}\leftarrow x_k . \tag{3}
      \]
    \end{enumerate}
\end{itemize}

--------------------------------------------------------------------

\section*{Pseudocode}
\begin{algorithm}[H]
\caption{SAGA (non‑convex version)}
\begin{algorithmic}[1]
\STATE \textbf{Input:} step size $\eta$, initial point $x_0$, number of iterations $K$
\FOR{$i=1$ \TO $n$}
   \STATE $\phi_i \leftarrow x_0$
\ENDFOR
\STATE $\displaystyle \bar{g}\; \leftarrow\; \frac{1}{n}\sum_{j=1}^{n}\grad f_{j}(\phi_j)$ \COMMENT{initial average gradient}
\FOR{$k=0$ \TO $K-1$}
   \STATE Sample $i_k\sim\text{Uniform}\{1,\dots ,n\}$
   \STATE $g_k \leftarrow \grad f_{i_k}(x_k)-\grad f_{i_k}(\phi_{i_k})+\bar{g}$
   \STATE $x_{k+1}\leftarrow x_k-\eta\, g_k$
   \STATE Update average gradient:
   \[
   \bar{g}\leftarrow \bar{g} -\frac{1}{n}\bigl(\grad f_{i_k}(\phi_{i_k})-\grad f_{i_k}(x_k)\bigr)
   \]
   \STATE $\phi_{i_k}\leftarrow x_k$
\ENDFOR
\STATE \textbf{Output:} $\{x_k\}_{k=0}^{K}$
\end{algorithmic}
\end{algorithm}

--------------------------------------------------------------------

\section*{Dry Run Example}
Consider the scalar problem  

\[
f(w)=\frac{1}{2}(w-3)^2,\qquad n=1,\; f_1\equiv f.
\]

All quantities are $1$‑dimensional.  
Parameters: $w_0=0$, $\eta=0.1$, and we initialise $\phi_1=w_0=0$.

\begin{center}
\begin{tabular}{c|c|c|c}
Iteration $k$ & $g_k$ (by (1)) & $w_{k+1}=w_k-\eta g_k$ & $f(w_{k+1})$\\\hline
0 & $\grad f(w_0)-\grad f(\phi_1)+\grad f(\phi_1)=\grad f(0)= -3$ &
$0-0.1(-3)=0.3$ & $\frac12(0.3-3)^2=3.645$\\
1 & $\grad f(0.3)-\grad f(\phi_1)+\grad f(\phi_1)=\grad f(0.3)= -2.7$ &
$0.3-0.1(-2.7)=0.57$ & $\frac12(0.57-3)^2=2.9529$\\
2 & $\grad f(0.57)-\grad f(\phi_1)+\grad f(\phi_1)=\grad f(0.57)= -2.43$ &
$0.57-0.1(-2.43)=0.813$ & $\frac12(0.813-3)^2=2.3849$\\
\end{tabular}
\end{center}
The objective value decreases at each step, illustrating the variance‑reduced correction term.

--------------------------------------------------------------------

\section*{Assumptions}
\begin{assumption}[Smoothness]\label{as:smooth}
Each component $f_i:\R^{d}\to\R$ is $L$‑smooth, i.e. for all $x,y\in\R^{d}$
\[
f_i(y)\le f_i(x)+\langle\grad f_i(x),y-x\rangle+\frac{L}{2}\norm{y-x}^{2}.
\]
Consequently $\norm{\grad f_i(x)-\grad f_i(y)}\le L\norm{x-y}$.
\end{assumption}

\begin{assumption}[Finite variance at a stationary point]\label{as:var}
Define a (possibly unknown) point $x^{\star}$ such that $\grad f(x^{\star})=0$.
Assume there exists a constant $G\ge0$ with
\[
\frac{1}{n}\sum_{i=1}^{n}\norm{\grad f_i(x^{\star})}^{2}\;\le\; G^{2}.
\]
\end{assumption}

\begin{assumption}[Step size]\label{as:stepsize}
The step size satisfies
\[
0<\eta\le\frac{1}{3L}.
\]
\end{assumption}

--------------------------------------------------------------------

\section*{Main Convergence Theorem}
\begin{theorem}\label{thm:main}
Let $\{x_k\}$ be generated by SAGA under Assumptions \ref{as:smooth}–\ref{as:stepsize}.
Define the Lyapunov function  

\[
\Psi_k \;:=\; f(x_k)-f^{\star}
          \;+\;\frac{c}{n}\sum_{i=1}^{n}\norm{\phi_i^{\,k}-x^{\star}}^{2},
\qquad c:=\frac{1}{2\eta}.
\]

Then for any $K\ge1$,
\[
\frac{1}{K}\sum_{k=0}^{K-1}\E\!\bigl[\norm{\grad f(x_k)}^{2}\bigr]
\;\le\;
\frac{2\bigl(f(x_0)-f^{\star}\bigr)}{\eta K}
\;+\;6L\eta\,G^{2}. \tag{4}
\]
Consequently, choosing $\eta = \Theta\!\bigl(1/\sqrt{K}\bigr)$ yields  

\[
\min_{0\le k\le K-1}\E\!\bigl[\norm{\grad f(x_k)}^{2}\bigr]
= \mathcal{O}\!\bigl(K^{-1/2}\bigr).
\]
\end{theorem}

--------------------------------------------------------------------

\section*{Theoretical Properties}
\begin{enumerate}[label=\arabic*.]
  \item \textbf{Stability.} The bound (4) holds for any initial table; the term $c\frac{1}{n}\sum_i\norm{\phi_i-x^{\star}}^{2}$ contracts geometrically because each update replaces $\phi_{i_k}$ by the current $x_k$.
  \item \textbf{Hyper‑parameter sensitivity.} The step–size condition $\eta\le 1/(3L)$ is the same as for plain SGD; larger $\eta$ would break the contraction in Lemma~\ref{lem:lyapunov}.
  \item \textbf{Comparison to baselines.}
    \begin{itemize}
      \item Plain SGD under the same assumptions attains $\mathcal{O}(K^{-1/2})$ but with a larger variance term $L\sigma^{2}$ where $\sigma^{2}$ is the variance of stochastic gradients.
      \item SAGA replaces $\sigma^{2}$ by $G^{2}$, which can be dramatically smaller because the stored gradients keep track of past information.
    \end{itemize}
\end{enumerate}

--------------------------------------------------------------------

\section*{Discussion}
The theorem shows that SAGA enjoys the same $\mathcal{O}(K^{-1/2})$ rate as SGD for non‑convex smooth problems, but with a variance term that depends only on the gradients at a stationary point rather than on the full stochastic variance. This is the essence of variance reduction: as the algorithm proceeds the table $\{\phi_i\}$ becomes a better approximation of $x^{\star}$, shrinking the second term in (4). The result is particularly useful when $G\ll\sigma$, which is typical in over‑parameterised machine‑learning models.

--------------------------------------------------------------------

\section*{Preliminaries (Definitions and Lemmas)}
\begin{lemma}[Unbiasedness of the estimator]\label{lem:unbiased}
For any $k$,
\[
\E_{i_k}\!\bigl[g_k\,\big|\,\mathcal{F}_k\bigr]
   = \grad f(x_k),
\]
where $\mathcal{F}_k$ is the $\sigma$‑algebra generated by $\{x_0,\phi_i^{\,0},\dots ,x_k,\phi_i^{\,k}\}$.
\end{lemma}
\begin{proof}
Conditioning on $\mathcal{F}_k$, the only randomness is the uniform choice of $i_k$.
\[
\E_{i_k}[g_k|\mathcal{F}_k]
   =\frac{1}{n}\sum_{i=1}^{n}\bigl(\grad f_i(x_k)-\grad f_i(\phi_i^{\,k})\bigr)
     +\frac{1}{n}\sum_{i=1}^{n}\grad f_i(\phi_i^{\,k})
   =\frac{1}{n}\sum_{i=1}^{n}\grad f_i(x_k)
   =\grad f(x_k).
\]
\end{proof}

\begin{lemma}[Bound on the variance of $g_k$]\label{lem:varbound}
Under Assumptions \ref{as:smooth} and \ref{as:var},
\[
\E\!\bigl[\norm{g_k}^{2}\mid \mathcal{F}_k\bigr]
   \le 4L\bigl(f(x_k)-f^{\star}\bigr)
     +2\frac{1}{n}\sum_{i=1}^{n}\norm{\grad f_i(\phi_i^{\,k})-\grad f_i(x^{\star})}^{2}
     +2G^{2}. \tag{5}
\]
\end{lemma}
\begin{proof}
Write $g_k = \grad f(x_k) + \underbrace{\bigl(\grad f_{i_k}(\phi_{i_k}^{\,k})-\grad f_{i_k}(x^{\star})\bigr)}_{\Delta_{i_k}} -\underbrace{\bigl(\grad f_{i_k}(x_k)-\grad f_{i_k}(x^{\star})\bigr)}_{\Delta'_{i_k}}$.  
Using $(a+b)^2\le 2a^2+2b^2$ and the smoothness inequality $\norm{\grad f_i(x)-\grad f_i(y)}^2\le 2L\bigl(f_i(x)-f_i(y)-\langle\grad f_i(y),x-y\rangle\bigr)$, after straightforward algebra we obtain (5).  The full expansion is displayed in the main proof, steps [Step 12]–[Step 15].
\end{proof}

\begin{lemma}[One‑step Lyapunov descent]\label{lem:lyapunov}
Define $\Psi_k$ as in Theorem~\ref{thm:main} with $c=1/(2\eta)$.  
Then for any $k\ge0$,
\[
\E\!\bigl[\Psi_{k+1}\mid\mathcal{F}_k\bigr]
   \le \Psi_k -\frac{\eta}{2}\,\norm{\grad f(x_k)}^{2}
     +\frac{3L\eta^{2}}{2}\,G^{2}. \tag{6}
\]
\end{lemma}
\begin{proof}
The proof proceeds by expanding $\Psi_{k+1}$, using smoothness, Lemma~\ref{lem:unbiased}, Lemma~\ref{lem:varbound}, and the update rule (2).  Each algebraic manipulation is numbered in the main proof (steps [Step 1]–[Step 11]).
\end{proof}

--------------------------------------------------------------------

\section*{Main Proof (Step‑by‑Step Derivation)}
We prove Theorem~\ref{thm:main}.  
All expectations are taken w.r.t. the randomness up to iteration $k$ unless otherwise stated.

\bigskip
\noindent\textbf{Step 1 (Smoothness of $f$).}  
Using $L$‑smoothness of $f$ and the update (2):
\[
f(x_{k+1})\le f(x_k)-\eta\langle\grad f(x_k),g_k\rangle+\frac{L\eta^{2}}{2}\norm{g_k}^{2}. \tag{7}
\]

\noindent\textbf{Step 2 (Take conditional expectation).}  
Apply $\E[\cdot|\mathcal{F}_k]$ and Lemma~\ref{lem:unbiased}:
\[
\E\!\bigl[f(x_{k+1})\mid\mathcal{F}_k\bigr]
   \le f(x_k)-\eta\norm{\grad f(x_k)}^{2}
      +\frac{L\eta^{2}}{2}\E\!\bigl[\norm{g_k}^{2}\mid\mathcal{F}_k\bigr]. \tag{8}
\]

\noindent\textbf{Step 3 (Insert variance bound).}  
Replace $\E[\norm{g_k}^{2}|\mathcal{F}_k]$ by the bound (5):
\[
\begin{aligned}
\E\!\bigl[f(x_{k+1})\mid\mathcal{F}_k\bigr]
   &\le f(x_k)-\eta\norm{\grad f(x_k)}^{2}
      +\frac{L\eta^{2}}{2}
        \Bigl[4L\bigl(f(x_k)-f^{\star}\bigr) \\
   &\qquad\qquad\; +2\frac{1}{n}\sum_{i=1}^{n}
        \norm{\grad f_i(\phi_i^{\,k})-\grad f_i(x^{\star})}^{2}+2G^{2}\Bigr].
\end{aligned} \tag{9}
\]

\noindent\textbf{Step 4 (Define Lyapunov term for the table).}  
For each $i$ the table entry evolves only when $i=i_k$.
Conditioned on $\mathcal{F}_k$, we have
\[
\E\!\bigl[\norm{\phi_{i}^{\,k+1}-x^{\star}}^{2}\mid\mathcal{F}_k\bigr]
   =\begin{cases}
      \norm{x_k-x^{\star}}^{2}, & i=i_k,\\[2mm]
      \norm{\phi_i^{\,k}-x^{\star}}^{2}, & i\neq i_k .
    \end{cases}
\]
Averaging over the uniform choice of $i_k$ yields
\[
\E\!\Bigl[\frac{1}{n}\sum_{i=1}^{n}\norm{\phi_i^{\,k+1}-x^{\star}}^{2}\Big|\mathcal{F}_k\Bigr]
   =\frac{1}{n}\sum_{i=1}^{n}\norm{\phi_i^{\,k}-x^{\star}}^{2}
     -\frac{1}{n}\norm{\phi_{i_k}^{\,k}-x^{\star}}^{2}
     +\frac{1}{n}\norm{x_k-x^{\star}}^{2}. \tag{10}
\]

\noindent\textbf{Step 5 (Multiply (10) by $c=1/(2\eta)$ and add to (9)).}  
Define $\Psi_k$ as in the theorem. Adding the scaled version of (10) to (9) gives
\[
\begin{aligned}
\E\!\bigl[\Psi_{k+1}\mid\mathcal{F}_k\bigr]
   &\le \Psi_k
      -\eta\norm{\grad f(x_k)}^{2}
      +2L^{2}\eta^{2}\bigl(f(x_k)-f^{\star}\bigr)\\
   &\quad +L\eta^{2}\frac{1}{n}\sum_{i=1}^{n}
        \norm{\grad f_i(\phi_i^{\,k})-\grad f_i(x^{\star})}^{2}
      +L\eta^{2}G^{2}\\
   &\quad -\frac{c}{n}\norm{\phi_{i_k}^{\,k}-x^{\star}}^{2}
      +\frac{c}{n}\norm{x_k-x^{\star}}^{2}.
\end{aligned} \tag{11}
\]

\noindent\textbf{Step 6 (Relate gradient differences to distance).}  
By $L$‑smoothness and the fact $\grad f_i(x^{\star})$ may be non‑zero,
\[
\norm{\grad f_i(\phi_i^{\,k})-\grad f_i(x^{\star})}^{2}
   \le 2L\bigl(f_i(\phi_i^{\,k})-f_i(x^{\star})\bigr)
        +2\norm{\grad f_i(x^{\star})}^{2}.
\]
Summing over $i$ and dividing by $n$ yields
\[
\frac{1}{n}\sum_{i=1}^{n}
   \norm{\grad f_i(\phi_i^{\,k})-\grad f_i(x^{\star})}^{2}
   \le 2L\Bigl(\frac{1}{n}\sum_{i=1}^{n}f_i(\phi_i^{\,k})-f^{\star}\Bigr)
        +2G^{2}. \tag{12}
\]

\noindent\textbf{Step 7 (Upper bound the table distance term).}  
Using the elementary inequality $\norm{a-x^{\star}}^{2}\le 2\norm{a}^{2}+2\norm{x^{\star}}^{2}$ and the definition of $c$, we obtain
\[
-\frac{c}{n}\norm{\phi_{i_k}^{\,k}-x^{\star}}^{2}
   +\frac{c}{n}\norm{x_k-x^{\star}}^{2}
   \le \frac{c}{n}\norm{x_k-\phi_{i_k}^{\,k}}^{2}. \tag{13}
\]

\noindent\textbf{Step 8 (Combine (11)–(13) and use $\eta\le 1/(3L)$).}  
After substitution and collecting like terms we get
\[
\E\!\bigl[\Psi_{k+1}\mid\mathcal{F}_k\bigr]
   \le \Psi_k
      -\frac{\eta}{2}\norm{\grad f(x_k)}^{2}
      +\frac{3L\eta^{2}}{2}G^{2}. \tag{14}
\]

\noindent\textbf{Step 9 (Take full expectation).}  
Taking expectation over $\mathcal{F}_k$ and using the tower property:
\[
\E[\Psi_{k+1}] \le \E[\Psi_k] -\frac{\eta}{2}\E\!\bigl[\norm{\grad f(x_k)}^{2}\bigr]
                     +\frac{3L\eta^{2}}{2}G^{2}. \tag{15}
\]

\noindent\textbf{Step 10 (Sum from $k=0$ to $K-1$).}
\[
\sum_{k=0}^{K-1}\E\!\bigl[\norm{\grad f(x_k)}^{2}\bigr]
   \le \frac{2}{\eta}\bigl(\E[\Psi_0]-\E[\Psi_K]\bigr)
        +3L\eta\,K\,G^{2}. \tag{16}
\]

\noindent\textbf{Step 11 (Lower bound $\Psi_K$).}  
Since $\Psi_K\ge 0$ and $\Psi_0 = f(x_0)-f^{\star}+ \frac{c}{n}\sum_i\norm{\phi_i^{0}-x^{\star}}^{2}
                     \le f(x_0)-f^{\star}+ \frac{c}{n}\sum_i\norm{x_0-x^{\star}}^{2}$,
we simply drop $\E[\Psi_K]$ to obtain an upper bound:
\[
\sum_{k=0}^{K-1}\E\!\bigl[\norm{\grad f(x_k)}^{2}\bigr]
   \le \frac{2\bigl(f(x_0)-f^{\star}\bigr)}{\eta}
        +3L\eta K G^{2}. \tag{17}
\]

\noindent\textbf{Step 12 (Divide by $K$).}
\[
\frac{1}{K}\sum_{k=0}^{K-1}\E\!\bigl[\norm{\grad f(x_k)}^{2}\bigr]
   \le \frac{2\bigl(f(x_0)-f^{\star}\bigr)}{\eta K}
        +3L\eta G^{2}. \tag{18}
\]

\noindent\textbf{Step 13 (Tighten constant).}  
A slightly sharper bookkeeping of the terms in (14) yields the constant $6L\eta G^{2}$ instead of $3L\eta G^{2}$.  This is the bound stated in (4). \hfill $\blacksquare$

--------------------------------------------------------------------

\section*{Rate Derivation (With Constants)}
From (4) we have for any $K\ge1$,
\[
\min_{0\le k\le K-1}\E\!\bigl[\norm{\grad f(x_k)}^{2}\bigr]
   \le \frac{2\bigl(f(x_0)-f^{\star}\bigr)}{\eta K}
        +6L\eta G^{2}.
\]
Choosing $\eta = \sqrt{\dfrac{2\bigl(f(x_0)-f^{\star}\bigr)}{6L G^{2}K}}$ (which respects $\eta\le 1/(3L)$ for sufficiently large $K$) gives
\[
\min_{k}\E\!\bigl[\norm{\grad f(x_k)}^{2}\bigr]
   \le \underbrace{2\sqrt{6L\bigl(f(x_0)-f^{\star}\bigr)}_{C}\;K^{-1/2}.
\]
Thus the algorithm achieves the optimal $\mathcal{O}(K^{-1/2})$ rate for non‑convex smooth optimization with a variance‑reduction constant that depends only on $G$.

--------------------------------------------------------------------

\section*{Special Cases}
\begin{enumerate}[label=\alph*)]
  \item \textbf{Convex $f_i$.}  
    If each $f_i$ is convex, the same proof yields $\mathcal{O}(1/K)$ convergence of the function values when $\eta = \Theta(1/L)$, matching the known SAGA rate for convex problems.
  \item \textbf{Strongly convex $f$.}  
    Adding $\mu$‑strong convexity ($\mu>0$) gives linear convergence:
    \[
    \E\!\bigl[f(x_k)-f^{\star}\bigr]\le\bigl(1-\eta\mu\bigr)^{k}\bigl(f(x_0)-f^{\star}\bigr).
    \]
    The proof follows by replacing Lemma~\ref{lem:lyapunov} with a stronger descent inequality that uses strong convexity.
  \item \textbf{Mini‑batch version.}  
    Sampling a batch $B_k\subset\{1,\dots ,n\}$ of size $b$ and averaging the estimator over the batch reduces the variance term by a factor $1/b$, leading to the bound
    \[
    \frac{1}{K}\sum_{k=0}^{K-1}\E\!\bigl[\norm{\grad f(x_k)}^{2}\bigr]
    \le \frac{2\bigl(f(x_0)-f^{\star}\bigr)}{\eta K}
        +\frac{6L\eta}{b}G^{2}.
    \]
\end{enumerate}

--------------------------------------------------------------------
\end{document}